\documentclass[10pt,a4paper]{amsart}
\usepackage[margin=19mm]{geometry}
\usepackage{amsmath,amssymb,mathtools}
\usepackage[scr=rsfs,cal=euler]{mathalfa}
\usepackage{xfrac}
\usepackage[unicode,hidelinks,bookmarks=false]{hyperref}
\newcommand{\ie}{i.e.\ }
\newcommand{\cf}{cf.\ }
\newcommand{\resp}{respectively\ }
\newcommand{\Subsection}{\subsection}
\providecommand{\nobreakdash}{\nobreak\hbox{-}}
\setlength{\emergencystretch}{2em}
\multlinegap0pt
\usepackage{algorithm}\usepackage{algpseudocode}
\newtheorem{theorem}{Theorem}
\newtheorem{lemma}{Lemma}
\title{Frequentist Regret Analysis of Gaussian Process Thompson Sampling via Fractional Posteriors}
\author{Somjit Roy, Prateek Jaiswal, Anirban Bhattacharya, Debdeep Pati, Bani K. Mallick}
\date{Source: arXiv:2602.14472v1 (CC BY 4.0)}
\begin{document}
\maketitle

\noindent\emph{Source and licence.} Frequentist Regret Analysis of Gaussian Process Thompson Sampling via Fractional Posteriors. Version arXiv:2602.14472v1 (CC BY 4.0), distributed by arXiv under the Creative Commons Attribution 4.0 International licence, http://creativecommons.org/licenses/by/4.0/, as recorded in the arXiv licence metadata for this exact version and shown on the abstract page https://arxiv.org/abs/2602.14472v1. Full licence notice: \texttt{licenses/arxiv-2602.14472-CC-BY-4.0.txt}.\par\smallskip

\noindent\textit{Editorial scope.} Counted unit: the article's lemma lem:identity (source0.tex lines 795--800). The article's complete original proof of it is delivered with it (source0.tex lines 802--876, one \texttt{proof} environment of 75 lines). No mathematics is written, repaired or re-derived: the statement and the proof are the article's own lines, in the article's own notation, and no retained line is substituted or re-flowed.  Retained uncounted as the source-local context the proof needs: lines 761--761; lines 775--781.  Nothing was omitted from any retained span, so the package carries no omission record.  No retained line contains a citation, so the wrapper carries no bibliography and no bibliography entry.  No retained line contains a figure environment, an image-inclusion command or drawing code, so no image is delivered and the package declares no asset dependency.  Editorial formatting only, in addition: an AMS \texttt{amsart} preamble that restates the article's own declarations and macros used by the retained lines, namely the environments \texttt{theorem}, \texttt{lemma} on separate counters, together with the standard packages those lines need.  The retained environments print in this document as Theorem 1, Lemma 1 in order of appearance, whereas the article prints the same items with the numbering of its own full document.  The complete native source of the article as distributed in the arXiv e-print for this exact version is bundled unchanged as \texttt{sources/194710/source0.tex}.
\section{Reproducing Kernel Hilbert Space (RKHS) and Mercer's Theorem}\label{app:RKHS-Mercer}

\begin{theorem}[Mercer's theorem]\label{thm:Mercer}
    Let $k(\cdot, \cdot)$ be a continuous kernel with respect to a finite Borel measure on $\mathcal X$. There exists $\{(\lambda_m, \phi_m)\}_{m=1}^{\infty}$ such that $\lambda_m\in \mathbb{R}^{+}$, $\phi_m\in \mathcal H_k$, for $m\geq 1$, and:
    \begin{align*}
        k(x, x') = \sum_{m=1}^{\infty}\lambda_m \phi_m(x)\phi_m(x'),
    \end{align*}
    where $\{\lambda_m\}_{m=1}^{\infty}$ and $\{\phi_m\}_{m=1}^{\infty}$ are referred to as the eigenvalues and the eigenfeatures (or eigenfunctions) of $k(\cdot, \cdot)$ respectively.
\end{theorem}

\begin{lemma}\label{lem:identity}
For any $f\in \mathcal{H}_t \subseteq \mathcal{H}_0$, it can be shown that:
$$
\|f\|_{k_{t}}^2= \|f\|_{k}^2 + \lambda^{-1}\sum_{s=1}^{t}(f(x_s))^2.
$$
\end{lemma}

\begin{proof}
Recall  that, $k_t(x,x')= k(x,x')- k(x,A_t)^\top (k(A_t,A_t)+\sigma^{-1} I_d)^{-1} k(A_t,x')$ and $\sigma^{-1}= \lambda$.  Since the kernel $k(s,t)$ is positive definite, continuous and square integrable, using Mercer's theorem (as given in Theorem \ref{thm:Mercer} in Appendix \ref{app:RKHS-Mercer}), there exists an orthonormal sequence of continuous eigenfunctions $\{\phi_j\}_{j=1}^{\infty}$ in $L^2(\mathcal X)$ with eigenvalues $\mu_1\geq \mu_2\ldots \geq 0$ and:
\begin{align}
    \int_{\mathcal X} k(s,t)\phi_j(s)ds= \mu_j \phi_j(t), \quad j={1,2,\ldots} \quad \text{and} \quad k(s,t)= \sum_{j=1}^{\infty} \mu_j \phi_j(s)\phi_j(t).
    \label{eq:rk0}
\end{align} 
The RKHS $\mathcal H_0$ spanned by kernel $k(\cdot,\cdot)$ consists of functions $g\in L^2(\mathcal X)$ satisfying $\sum_{j=1}^{\infty}{g_j^2}/{\mu_j} < \infty$, where $g_j= \langle g, \phi_j\rangle_{L^2(\mathcal X)}$. Moreover, for any $g\in\mathcal H_0$,  $\|g\|_{\mathcal H_0}^2 = \sum_{j=1}^{\infty}{g_j^2}/{\mu_j}$.

Now let us assume that $\{\gamma_i,\psi_i\}_{i=1}^{\infty}$ are the corresponding eigen system of the reproducing kernel $k_t(\cdot,\cdot)$. It follows that:
\begin{align}
    \int_{\mathcal X} k_t(x,x')\psi_j(x)dx= \gamma_j \psi_j(x'), \quad j={1,2,\ldots}
    \label{eq:rk1}
\end{align}

Considering the left hand side above and substituting the expression of $k_t$, it follows that:
\begin{align}
    \int_{\mathcal X} k_t(s,t)\psi_j(s)ds = \int_{\mathcal X} \left[k(x,x')- k(x,A_t)^\top (k(A_t,A_t)+\sigma^{-1} I_d)^{-1} k(A_t,x')\right]\psi_j(s)ds.
    \label{eq:rk2}
\end{align}

Since, $\psi_j \in L^2(\mathcal X)$ and $\{\phi_i\}_{i=1}^{\infty}$ is an orthonormal basis of $L^2(\mathcal X)$, we can express $\psi_j(\cdot) = \sum_{i=1}^{\infty} p^{(j)}_i  \phi_i(\cdot) $. Substituting this expression into~\eqref{eq:rk2} and interchanging integral and summation (assuming sufficient regularity of RKHS), we obtain from using the first expression in~\eqref{eq:rk0} that:
\begin{align}\label{eq:rk3}
\begin{split}
    &\int_{\mathcal X} k_t(x,x')\psi_j(x)dx\\
    &\qquad = \int_{\mathcal X} \left[k(x,x')- k(x,A_t)^\top (k(A_t,A_t)+\sigma^{-1} I_d)^{-1} k(A_t,x')\right]\sum_{i=1}^{\infty} p^{(j)}_i  \phi_i(x)dx
    \\
    &\qquad = \sum_{i=1}^{\infty} p^{(j)}_i \mu_i \phi_i(x') - k(x', A_t) (k(A_t,A_t)+\sigma^{-1} I_d)^{-1} \sum_{i=1}^{\infty} p^{(j)}_i \mu_i \phi_i(A_t)
    \\
    &\qquad = \sum_{i=1}^{\infty} p^{(j)}_i \mu_i \phi_i(x') - k(x', A_t) (k(A_t,A_t)+\sigma^{-1} I_d)^{-1} \sum_{i=1}^{\infty} p^{(j)}_i \mu_i \phi_i(A_t).
\end{split}
\end{align}

Now substituting~\eqref{eq:rk3} into~\eqref{eq:rk1} and taking inner product with $\phi_k(x')$, we have for all $k=1,2, \ldots$:
\begin{align}\label{eq:rk4}
\begin{split}
    &p^{(j)}_k \mu_k  - \mu_k \phi_k(A_t)^\top (k(A_t,A_t)+\sigma^{-1} I_d)^{-1} \sum_{i=1}^{\infty} p^{(j)}_i \mu_i \phi_i(A_t)\\
    &\qquad \qquad \qquad = \gamma_j \langle\psi_j(x'),\phi_k(x')\rangle = \gamma_j p^{(j)}_k.
\end{split}
\end{align}

We use the following matrix notation:
\begin{align}
\nonumber
    \mathbf{M}=\text{diag}(\mu_1,\mu_2,\ldots), \quad
        \mathbf{p}^{(j)}= (p^{j}_1,p^{j}_2,\ldots)^\top, \quad
    \Phi(A_t)= (\phi_1(A_t),\phi_2(A_t),\ldots)^\top.
\end{align}

Using the second expression in~\eqref{eq:rk0}, observe that $k(A_t,A_t)=  \Phi(A_t)^\top \mathbf{M} \Phi(A_t)$. Combined with this and using the above matrix notation,~\eqref{eq:rk4} can be rewritten as:
\begin{align}
    \left(\mathbf{M}  - \mathbf{M} \Phi(A_t)(\Phi(A_t)^\top \mathbf{M} \Phi(A_t)+\sigma^{-1} I_d)^{-1}\Phi(A_t)^\top \mathbf{M} \right) \mathbf{p}^{(j)} = \gamma_j \mathbf{p}^{(j)}.
\end{align}

With an application of Woodbury matrix identity, we have: 
$$\left(\mathbf{M}  - \mathbf{M} \Phi(A_t)(\Phi(A_t)^\top \mathbf{M} \Phi(A_t)+\sigma^{-1} I_d)^{-1}\Phi(A_t)^\top \mathbf{M} \right) = (\mathbf{M}^{-1} + \sigma \Phi(A_t)\Phi(A_t)^\top)^{-1},$$ 
and therefore:
\begin{align}
    \gamma_j \mathbf{p}^{(j)} = (\mathbf{M}^{-1} + \sigma \Phi(A_t)\Phi(A_t)^\top)^{-1} \mathbf{p}^{(j)}, \quad j=1,2,\ldots.
    \label{eq:rk5}
\end{align}

For any $f\in \mathcal H_t \subseteq\mathcal H_0$, observe that:
\begin{align}\label{eq:rk6}
\begin{split}
    \|f\|_{\mathcal H_t}^2 &=\sum_{j=1}^{\infty}  \frac{\langle f,\psi_j\rangle_{L^2(\mathcal X)}^2}{\gamma_j} = \sum_{j=1}^{\infty}  \frac{\langle f,\sum_{i=1}^{\infty} p^{(j)}_i  \phi_i \rangle_{L^2(\mathcal X)}^2}{\gamma_j}\\
    &= \sum_{j=1}^{\infty}  \frac{(\sum_{i=1}^{\infty} p^{(j)}_i f_i )^2}{\gamma_j}  = \sum_{j=1}^{\infty}  \frac{((\mathbf{p}^{(j)})^\top \mathbf{f})^2}{\gamma_j}\\
    & =   \mathbf{f}^\top \mathbf{P}\Gamma^{-1}\mathbf{P}^\top\mathbf{f},
\end{split}
\end{align}
where $\mathbf{f}=(f_1,f_2,\ldots)^\top$, $\mathbf{P}= (\mathbf{p}^{(1)},\mathbf{p}^{(2)},\ldots)^\top$, and $\Gamma= \text{diag}(\gamma_1,\gamma_2,\ldots)$. Using these notations, \eqref{eq:rk5} can be rewritten as, $(\mathbf{M}^{-1} + \sigma \Phi(A_t)\Phi(A_t)^\top)^{-1} \mathbf{P}=  \mathbf{P}\Gamma$. Also note using \eqref{eq:rk5} that, matrix $\mathbf{P}$ is an orthogonal matrix, which implies that $\mathbf{P}\Gamma^{-1}\mathbf{P}^\top = (\mathbf{M}^{-1} + \sigma \Phi(A_t)\Phi(A_t)^\top) $. Substituting the expression for $\mathbf{P}\Gamma^{-1}\mathbf{P}^\top$ in~\eqref{eq:rk6}, it follows that:
\begin{align}
    \|f\|_{\mathcal H_t}^2 = \mathbf{f}^\top (\mathbf{M}^{-1} + \sigma \Phi(A_t)\Phi(A_t)^\top) \mathbf{f} = \|f\|_{\mathcal H_0}^2 + \sigma \mathbf{f}^\top \Phi(A_t)\Phi(A_t)^\top \mathbf{f}.
\end{align}
Hence the assertion of Lemma \ref{lem:identity} follows, since $\sigma=  \lambda^{-1}$ and $\Phi(A_t)^\top \mathbf{f} = \sum_{j=1}^{\infty} f_j \phi_j(A_t) = f(A_t)$ because $\{\phi_j\}_{j=1}^{\infty}$ are orthonormal basis of $L^2(\mathcal X)$, which implies that $f(A_t)^\top f(A_t)= \sum_{s=1}^{t}(f(x_t))^2$.
\end{proof}

\end{document}
